% ============================================================
\chapter{Mod\`ele de Black--Scholes}
\label{ch:black_scholes}
% ============================================================

En 1973, Fischer Black et Myron Scholes publient un article qui va
transformer la finance~: ils d\'erivent une formule explicite pour le prix
d'une option europ\'eenne. Robert Merton, ind\'ependamment, \'etend et
solidifie la th\'eorie. Le mod\`ele repose sur une hypoth\`ese audacieuse~:
le cours de l'action suit un mouvement brownien g\'eom\'etrique, et dans ce
cadre, il est possible de construire un portefeuille de couverture qui
\'elimine \emph{tout} le risque. Le prix de l'option est alors le co\^ut de
cette couverture. Scholes et Merton recevront le prix Nobel d'\'economie en
1997 (Black \'etant d\'ec\'ed\'e en 1995).

% ------------------------------------------------------------
\section{Introduction}
% ------------------------------------------------------------

Le mod\`ele de Black--Scholes--Merton (1973) constitue la pierre angulaire
de la finance quantitative moderne. Il fournit un cadre math\'ematique
rigoureux pour l'\'evaluation et la couverture des options europ\'eennes,
bas\'e sur l'hypoth\`ese que le prix de l'actif sous-jacent suit un
mouvement brownien g\'eom\'etrique.

\section{Mouvement brownien g\'eom\'etrique}

\begin{definition}[Mouvement brownien g\'eom\'etrique (MBG)]
Un processus $(S_t)_{t \geq 0}$ suit un \textbf{mouvement brownien
g\'eom\'etrique} s'il satisfait l'\'equation diff\'erentielle stochastique~:
\[
    dS_t = \mu S_t \, dt + \sigma S_t \, dW_t,
\]
o\`u $\mu \in \R$ est le rendement esp\'er\'e (drift), $\sigma > 0$ la
volatilit\'e, et $(W_t)_{t \geq 0}$ un mouvement brownien standard.
\end{definition}

\begin{theoreme}[Solution du MBG]
\label{thm:solution_mbg}
La solution de l'EDS ci-dessus est~:
\[
    S_t = S_0 \exp\!\left[\left(\mu - \frac{\sigma^2}{2}\right)t
          + \sigma W_t\right].
\]
En particulier, $\ln(S_t/S_0) \sim \mathcal{N}\!\left((\mu - \sigma^2/2)t,\;
\sigma^2 t\right)$.
\end{theoreme}

\begin{proof}
Posons $X_t = \ln S_t$. Par la formule d'It\^o (cf.\ chapitre 5)~:
\[
    dX_t = \frac{1}{S_t}dS_t - \frac{1}{2}\frac{1}{S_t^2}(dS_t)^2
         = \left(\mu - \frac{\sigma^2}{2}\right)dt + \sigma\, dW_t.
\]
En int\'egrant de $0$ \`a $t$~:
$X_t = X_0 + (\mu - \sigma^2/2)t + \sigma W_t$,
d'o\`u le r\'esultat.
\end{proof}

\section{Hypoth\`eses du mod\`ele de Black--Scholes}

\begin{attention}[title=\textbf{Hypoth\`eses fondamentales}]
\begin{enumerate}
    \item Le sous-jacent suit un MBG de volatilit\'e $\sigma$ constante.
    \item Le taux sans risque $r$ est constant.
    \item Pas de dividendes (g\'en\'eralisation possible).
    \item March\'e sans frictions (pas de co\^uts de transaction).
    \item Trading continu possible.
    \item Vente \`a d\'ecouvert autoris\'ee.
\end{enumerate}
\end{attention}

\section{\'Equation de Black--Scholes}

\begin{theoreme}[\'EDP de Black--Scholes]
\label{thm:edp_bs}
Sous les hypoth\`eses du mod\`ele, le prix $V(S,t)$ d'un d\'eriv\'e
europ\'een satisfait l'\'equation aux d\'eriv\'ees partielles~:
\[
    \frac{\partial V}{\partial t} + rS\frac{\partial V}{\partial S}
    + \frac{1}{2}\sigma^2 S^2 \frac{\partial^2 V}{\partial S^2} - rV = 0,
\]
avec la condition terminale $V(S,T) = h(S)$ o\`u $h$ est le payoff.
\end{theoreme}

\begin{proof}[Esquisse de preuve par r\'eplication]
Consid\'erons un portefeuille $\Pi = V - \Delta S$ o\`u $\Delta =
\partial V/\partial S$. Par It\^o~:
\begin{align*}
    d\Pi &= dV - \Delta \, dS \\
         &= \left(\frac{\partial V}{\partial t}
            + \frac{1}{2}\sigma^2 S^2 \frac{\partial^2 V}{\partial S^2}\right)dt.
\end{align*}
Ce portefeuille est sans risque (pas de terme en $dW_t$), donc par AOA~:
$d\Pi = r\Pi\,dt = r(V - \Delta S)\,dt$.
En \'egalisant, on obtient l'EDP.
\end{proof}

\section{Formules de Black--Scholes}

\begin{theoreme}[Formules ferm\'ees de Black--Scholes]
\label{thm:formules_bs}
Le prix d'un call europ\'een est~:
\[
    C(S,t) = S\,\Phi(d_1) - Ke^{-r(T-t)}\Phi(d_2),
\]
et le prix d'un put europ\'een~:
\[
    P(S,t) = Ke^{-r(T-t)}\Phi(-d_2) - S\,\Phi(-d_1),
\]
o\`u $\Phi$ d\'esigne la fonction de r\'epartition de la loi normale
standard et~:
\[
    d_1 = \frac{\ln(S/K) + (r + \sigma^2/2)(T-t)}{\sigma\sqrt{T-t}}, \quad
    d_2 = d_1 - \sigma\sqrt{T-t}.
\]
\end{theoreme}

\begin{formules}[title=\textbf{Formules de Black--Scholes --- R\'ecapitulatif}]
\begin{align*}
    C &= S\Phi(d_1) - Ke^{-r\tau}\Phi(d_2), \\
    P &= Ke^{-r\tau}\Phi(-d_2) - S\Phi(-d_1), \\
    d_1 &= \frac{\ln(S/K) + (r + \sigma^2/2)\tau}{\sigma\sqrt{\tau}}, \quad
    d_2 = d_1 - \sigma\sqrt{\tau}, \quad \tau = T - t.
\end{align*}
\end{formules}

\begin{proof}[Preuve par pricing risque-neutre]
Sous la mesure risque-neutre $\mathbb{Q}$~:
\[
    C = e^{-r\tau}\E^\mathbb{Q}[(S_T - K)^+].
\]
Or $S_T = S\exp\!\left[(r - \sigma^2/2)\tau + \sigma\sqrt{\tau}\,Z\right]$
avec $Z \sim \mathcal{N}(0,1)$ sous $\mathbb{Q}$. Le payoff est positif
lorsque $Z > -d_2$. On calcule~:
\begin{align*}
    C &= e^{-r\tau}\int_{-d_2}^{+\infty}
         \left(Se^{(r-\sigma^2/2)\tau + \sigma\sqrt{\tau}\,z} - K\right)
         \frac{e^{-z^2/2}}{\sqrt{2\pi}}\,dz \\
      &= S\Phi(d_1) - Ke^{-r\tau}\Phi(d_2).
\end{align*}
Le passage de $d_2$ \`a $d_1$ dans le premier terme r\'esulte du
compl\'etement du carr\'e dans l'exponentielle.
\end{proof}

\section{Les Grecques}

Les \textbf{Grecques} (Greeks) mesurent la sensibilit\'e du prix de
l'option aux diff\'erents param\`etres du mod\`ele.

\begin{definition}[Grecques principales]
Pour un call europ\'een~:
\begin{align*}
    \Delta &= \frac{\partial C}{\partial S} = \Phi(d_1), \\
    \Gamma &= \frac{\partial^2 C}{\partial S^2}
            = \frac{\phi(d_1)}{S\sigma\sqrt{\tau}}, \\
    \Theta &= \frac{\partial C}{\partial t}
            = -\frac{S\sigma\phi(d_1)}{2\sqrt{\tau}} - rKe^{-r\tau}\Phi(d_2), \\
    \mathcal{V} &= \frac{\partial C}{\partial \sigma}
                  = S\sqrt{\tau}\,\phi(d_1), \\
    \rho &= \frac{\partial C}{\partial r} = K\tau e^{-r\tau}\Phi(d_2),
\end{align*}
o\`u $\phi(x) = \frac{1}{\sqrt{2\pi}}e^{-x^2/2}$ est la densit\'e
de la loi normale standard.
\end{definition}

\begin{intuition}[title=\textbf{Interpr\'etation des Grecques}]
\begin{itemize}
    \item $\Delta$~: nombre d'unit\'es du sous-jacent pour couvrir l'option.
    \item $\Gamma$~: convexit\'e, mesure la stabilit\'e du delta.
    \item $\Theta$~: d\'ecroissance temporelle (\textit{time decay}).
    \item $\mathcal{V}$ (vega)~: sensibilit\'e \`a la volatilit\'e.
    \item $\rho$~: sensibilit\'e au taux d'int\'er\^et.
\end{itemize}
\end{intuition}

\begin{proposition}[Relation fondamentale des Grecques]
Le delta, gamma et theta d'un d\'eriv\'e satisfont l'EDP de Black--Scholes
sous la forme~:
\[
    \Theta + rS\Delta + \frac{1}{2}\sigma^2 S^2 \Gamma = rV.
\]
\end{proposition}

\section{Volatilit\'e implicite}

\begin{definition}[Volatilit\'e implicite]
La \textbf{volatilit\'e implicite} $\sigma_{\text{impl}}$ est la valeur
de $\sigma$ qui, ins\'er\'ee dans la formule de Black--Scholes, reproduit
le prix de march\'e observ\'e $C_{\text{mkt}}$~:
\[
    C_{\text{BS}}(S, K, T, r, \sigma_{\text{impl}}) = C_{\text{mkt}}.
\]
\end{definition}

\begin{remarque}
L'existence et l'unicit\'e de $\sigma_{\text{impl}}$ d\'ecoulent du fait
que $\mathcal{V} = \partial C/\partial\sigma > 0$ (le prix du call est
strictement croissant en la volatilit\'e).
\end{remarque}

\section{Impl\'ementation Python}

\begin{codeblock}[title=\textbf{Formules de Black--Scholes et Grecques}]
\begin{minted}{python}
import numpy as np
from scipy.stats import norm

def bs_price(S, K, T, r, sigma, option='call'):
    """Prix Black-Scholes d'une option europeenne."""
    d1 = (np.log(S/K) + (r + 0.5*sigma**2)*T) / (sigma*np.sqrt(T))
    d2 = d1 - sigma*np.sqrt(T)
    if option == 'call':
        return S*norm.cdf(d1) - K*np.exp(-r*T)*norm.cdf(d2)
    else:
        return K*np.exp(-r*T)*norm.cdf(-d2) - S*norm.cdf(-d1)

def bs_greeks(S, K, T, r, sigma):
    """Calcul des Grecques pour un call europeen."""
    d1 = (np.log(S/K) + (r + 0.5*sigma**2)*T) / (sigma*np.sqrt(T))
    d2 = d1 - sigma*np.sqrt(T)
    delta = norm.cdf(d1)
    gamma = norm.pdf(d1) / (S * sigma * np.sqrt(T))
    theta = (-S*sigma*norm.pdf(d1)/(2*np.sqrt(T))
             - r*K*np.exp(-r*T)*norm.cdf(d2))
    vega  = S * np.sqrt(T) * norm.pdf(d1)
    rho   = K * T * np.exp(-r*T) * norm.cdf(d2)
    return {'delta': delta, 'gamma': gamma, 'theta': theta,
            'vega': vega, 'rho': rho}

# Exemple
S, K, T, r, sigma = 100, 100, 1.0, 0.05, 0.20
print(f"Call = {bs_price(S,K,T,r,sigma,'call'):.4f}")
print(f"Put  = {bs_price(S,K,T,r,sigma,'put'):.4f}")
greeks = bs_greeks(S, K, T, r, sigma)
for name, val in greeks.items():
    print(f"{name:>6s} = {val:.6f}")
\end{minted}
\end{codeblock}

\begin{codeblock}[title=\textbf{Calcul de la volatilit\'e implicite (Newton-Raphson)}]
\begin{minted}{python}
def implied_vol(market_price, S, K, T, r, option='call',
                tol=1e-8, max_iter=100):
    """Volatilit\'e implicite par Newton-Raphson."""
    sigma = 0.20  # initial guess
    for _ in range(max_iter):
        price = bs_price(S, K, T, r, sigma, option)
        d1 = (np.log(S/K) + (r + 0.5*sigma**2)*T) / (sigma*np.sqrt(T))
        vega = S * np.sqrt(T) * norm.pdf(d1)
        if vega < 1e-12:
            break
        sigma -= (price - market_price) / vega
        if abs(price - market_price) < tol:
            return sigma
    return sigma

# Exemple: retrouver sigma a partir d'un prix
C_mkt = bs_price(100, 100, 1.0, 0.05, 0.25, 'call')
iv = implied_vol(C_mkt, 100, 100, 1.0, 0.05)
print(f"Vol implicite retrouvee: {iv:.6f}")
\end{minted}
\end{codeblock}

\section{Extensions du mod\`ele}

\subsection{Dividendes continus}

Pour un sous-jacent versant un taux de dividende continu $q$, la formule
de Black--Scholes se modifie~:
\[
    C = Se^{-q\tau}\Phi(d_1) - Ke^{-r\tau}\Phi(d_2),
\]
avec $d_1 = \frac{\ln(S/K) + (r - q + \sigma^2/2)\tau}{\sigma\sqrt{\tau}}$.

\subsection{Options sur futures (formule de Black)}

\begin{theoreme}[Formule de Black (1976)]
Pour une option europ\'eenne sur un contrat futures de prix $F$~:
\[
    C = e^{-r\tau}\bigl[F\Phi(d_1) - K\Phi(d_2)\bigr],
\]
avec $d_1 = \frac{\ln(F/K) + \sigma^2\tau/2}{\sigma\sqrt{\tau}}$,
$d_2 = d_1 - \sigma\sqrt{\tau}$.
\end{theoreme}

\section{Limites du mod\`ele}

\begin{attention}[title=\textbf{Limites connues du mod\`ele de Black--Scholes}]
\begin{enumerate}
    \item La volatilit\'e n'est pas constante en pratique
          (smile de volatilit\'e, cf.\ chapitre 6).
    \item Les rendements ne sont pas gaussiens (queues \'epaisses).
    \item Le trading continu est une id\'ealisation.
    \item Les co\^uts de transaction ne sont pas nuls.
    \item Le mod\`ele sous-estime le risque de crash.
\end{enumerate}
\end{attention}

\section{Exercices}

\begin{exercice}[V\'erification de la formule de Black--Scholes]
\begin{enumerate}
    \item V\'erifier par substitution directe que la formule du call
          satisfait l'EDP de Black--Scholes.
    \item V\'erifier la condition terminale $C(S,T) = (S-K)^+$.
    \item D\'emontrer la parit\'e call-put \`a partir des formules.
\end{enumerate}
\end{exercice}

\begin{exercice}[Comportements limites]
\'Etudier les limites suivantes~:
\begin{enumerate}
    \item $C(S,t)$ lorsque $\sigma \to 0^+$.
    \item $C(S,t)$ lorsque $\sigma \to +\infty$.
    \item $C(S,t)$ lorsque $T - t \to 0^+$ avec $S > K$, $S = K$, $S < K$.
    \item $\Delta$ lorsque $T - t \to 0^+$.
\end{enumerate}
\end{exercice}

\begin{exercice}[Grecques num\'eriques]
\begin{enumerate}
    \item Calculer le delta et le gamma par diff\'erences finies et comparer
          aux formules analytiques.
    \item Tracer $\Delta$, $\Gamma$, $\Theta$, $\mathcal{V}$ en fonction
          de $S$ pour diff\'erentes valeurs de $\tau$.
    \item V\'erifier num\'eriquement la relation $\Theta + rS\Delta +
          \frac{1}{2}\sigma^2 S^2\Gamma = rC$.
\end{enumerate}
\end{exercice}

\begin{exercice}[Surface de volatilit\'e implicite]
\begin{enumerate}
    \item G\'en\'erer des prix d'options \`a partir d'un mod\`ele de Heston
          (chapitre 6) pour diff\'erents strikes et maturit\'es.
    \item Calculer la volatilit\'e implicite Black--Scholes pour chaque prix.
    \item Tracer la surface de volatilit\'e implicite $(K, T) \mapsto
          \sigma_{\text{impl}}(K,T)$.
\end{enumerate}
\end{exercice}
